\documentclass{standalone}
\usepackage{circuitikz}
\usetikzlibrary{calc}
\definecolor{circuitnotemark}{HTML}{FE00FE}
\begin{document}
\begin{circuitikz}[american, line width=0.8pt]
\ctikzset{bipoles/length=1.2cm}
\foreach \x in {0,1.2,2.4,3.6,4.8,6,7.2,8.4} {\foreach \y in {0,-1.2,-2.4,-3.6,-4.8,-6,-7.2} {\fill[gray, opacity=0.35] (\x,\y) circle (0.66pt);}}
\node[gray, font=\scriptsize] at (0,0.504) {1};
\node[gray, font=\scriptsize] at (1.2,0.504) {2};
\node[gray, font=\scriptsize] at (2.4,0.504) {3};
\node[gray, font=\scriptsize] at (3.6,0.504) {4};
\node[gray, font=\scriptsize] at (4.8,0.504) {5};
\node[gray, font=\scriptsize] at (6,0.504) {6};
\node[gray, font=\scriptsize] at (7.2,0.504) {7};
\node[gray, font=\scriptsize] at (8.4,0.504) {8};
\node[gray, font=\scriptsize, anchor=east] at (-0.504,0) {a};
\node[gray, font=\scriptsize, anchor=east] at (-0.504,-1.2) {b};
\node[gray, font=\scriptsize, anchor=east] at (-0.504,-2.4) {c};
\node[gray, font=\scriptsize, anchor=east] at (-0.504,-3.6) {d};
\node[gray, font=\scriptsize, anchor=east] at (-0.504,-4.8) {e};
\node[gray, font=\scriptsize, anchor=east] at (-0.504,-6) {f};
\node[gray, font=\scriptsize, anchor=east] at (-0.504,-7.2) {g};
\coordinate (b2) at (1.2,-1.2);
\coordinate (b5) at (4.8,-1.2);
\coordinate (b8) at (8.4,-1.2);
\coordinate (f2) at (1.2,-6);
\coordinate (f5) at (4.8,-6);
\coordinate (f8) at (8.4,-6);
\coordinate (a2) at (1.2,0);
\coordinate (c2) at (1.2,-2.4);
\coordinate (b1) at (0,-1.2);
\coordinate (a5) at (4.8,0);
\coordinate (c5) at (4.8,-2.4);
\coordinate (b4) at (3.6,-1.2);
\coordinate (a8) at (8.4,0);
\coordinate (c8) at (8.4,-2.4);
\coordinate (b7) at (7.2,-1.2);
\coordinate (e2) at (1.2,-4.8);
\coordinate (g2) at (1.2,-7.2);
\coordinate (f1) at (0,-6);
\coordinate (e5) at (4.8,-4.8);
\coordinate (g5) at (4.8,-7.2);
\coordinate (f4) at (3.6,-6);
\coordinate (e8) at (8.4,-4.8);
\coordinate (g8) at (8.4,-7.2);
\coordinate (f7) at (7.2,-6);
\node[npn] (part-Q1) at (b2) {}; % line 3
\node[pnp] (part-Q2) at (b5) {}; % line 4
\node[nigbt] (part-Q3) at (b8) {}; % line 5
\node[nmos] (part-M1) at (f2) {}; % line 6
\node[pmos] (part-M2) at (f5) {}; % line 7
\node[pigbt] (part-Q4) at (f8) {}; % line 8
\draw (a2) -| (part-Q1.collector); % line 10
\draw (c2) -| (part-Q1.emitter); % line 11
\draw (b1) -| (part-Q1.base); % line 12
\draw (a5) -| (part-Q2.collector); % line 13
\draw (c5) -| (part-Q2.emitter); % line 14
\draw (b4) -| (part-Q2.base); % line 15
\draw (a8) -| (part-Q3.collector); % line 16
\draw (c8) -| (part-Q3.emitter); % line 17
\draw (b7) -| (part-Q3.gate); % line 18
\draw (e2) -| (part-M1.drain); % line 19
\draw (g2) -| (part-M1.source); % line 20
\draw (f1) -| (part-M1.gate); % line 21
\draw (e5) -| (part-M2.drain); % line 22
\draw (g5) -| (part-M2.source); % line 23
\draw (f4) -| (part-M2.gate); % line 24
\draw (e8) -| (part-Q4.collector); % line 25
\draw (g8) -| (part-Q4.emitter); % line 26
\draw (f7) -| (part-Q4.gate); % line 27
\path (-0.07,-4.193) rectangle (2.47,-3.007); % line 29
\node[anchor=west, inner sep=0, circuitnotemark, font=\large] at (1.2,-3.6) {X}; % line 29
\path (3.53,-4.193) rectangle (6.07,-3.007); % line 30
\node[anchor=west, inner sep=0, circuitnotemark, font=\large] at (4.8,-3.6) {X}; % line 30
\path (6.876,-4.193) rectangle (9.924,-3.007); % line 31
\node[anchor=west, inner sep=0, circuitnotemark, font=\large] at (8.4,-3.6) {X}; % line 31
\path (-0.197,-8.993) rectangle (2.597,-7.807); % line 32
\node[anchor=west, inner sep=0, circuitnotemark, font=\large] at (1.2,-8.4) {X}; % line 32
\path (3.403,-8.993) rectangle (6.197,-7.807); % line 33
\node[anchor=west, inner sep=0, circuitnotemark, font=\large] at (4.8,-8.4) {X}; % line 33
\path (6.876,-8.993) rectangle (9.924,-7.807); % line 34
\node[anchor=west, inner sep=0, circuitnotemark, font=\large] at (8.4,-8.4) {X}; % line 34
\node[anchor=south west, inner sep=2pt, circuitnotemark, font=\large\bfseries] (circuittitle) at (current bounding box.north west) {X};
\path (circuittitle.south west) rectangle ++(2.667,0.593);
\end{circuitikz}
\end{document}
